% =====================================================================
%  06  反応器設計 — 左右2カラム
%  左：性質→部品マップ（上）＋形式選定3判断の表（下）
%  右：反応器の説明文（上・小さめ）＋浅床軸流・多基切替式固定床の模式図（下）
%  source_memo: report/chapters/04_reaction.tex §4.1 形式選定3判断 /
%               §4.4 接触時間 τ∝1/u² / 反応器概念断面 TikZ。
%  方針：[t] で帯直下から少し下げる。丸括弧禁止。はみ出し・重なりなし。
% =====================================================================
\begin{frame}[t]

  \vspace{0.35em}
  \begin{columns}[T,onlytextwidth]
    % ===== 左：性質→部品マップ（上）＋形式選定の3判断の表（下）=====
    \begin{column}{0.40\textwidth}
      \centering
      \resizebox{\linewidth}{!}{%
      \begin{tikzpicture}[
        >=Stealth, font=\small,
        prop/.style={draw=black!50, rounded corners=2pt, align=center,
                     inner sep=3pt, minimum height=0.7cm, text width=2.0cm},
        part/.style={draw=HeaderBlue, fill=HiBlue!45, rounded corners=2pt, align=center,
                     inner sep=3pt, minimum height=0.6cm, text width=2.9cm,
                     font=\small\bfseries},
        link/.style={-{Stealth[length=2mm]}, ArrowBlue, line width=1pt},
      ]
        \node[prop] (heat) at (0,2.85) {吸熱};
        \node[prop] (eq)   at (0,1.42){平衡支配\\$0.5$ bar};
        \node[prop] (coke) at (0,0)   {コーキング\\数分で失活};
        \node[part] (hgm)   at (3.7,2.85) {床内 HGM 補温};
        \node[part] (shal)  at (3.7,1.9)  {浅床軸流};
        \node[part] (multi) at (3.7,0.95) {多基・大断面・低速};
        \node[part] (swing) at (3.7,0.0)  {切替式　反応 $\rightleftarrows$ 再生};
        \draw[link] (heat.east) -- (hgm.west);
        \draw[link] (eq.east)   -- (shal.west);
        \draw[link] (eq.east)   -- (multi.west);
        \draw[link] (coke.east) -- (swing.west);
      \end{tikzpicture}}\par

      \vspace{1.0em}
      \raggedright
      {\small\bfseries 形式選定の3判断と採用根拠}\par
      \vspace{0.4em}
      {\scriptsize\renewcommand{\arraystretch}{1.15}%
       \begin{tabular}{@{}l l c@{}}
         \toprule
         判断 & 候補 & 採否 \\
         \midrule
         \multirow{2}{*}{触媒再生}
            & 移動床            & $\times$ \\
            & {\color{SpRed}並列切替固定床} & $\bigcirc$ \\
         \midrule
         \multirow{3}{*}{流通形式}
            & 軸流深床   & $\times$ \\
            & 径方向流   & $\triangle$ \\
            & {\color{SpRed}浅床軸流}   & $\bigcirc$ \\
         \midrule
         \multirow{2}{*}{熱補償}
            & 段間再加熱 & $\times$ \\
            & {\color{SpRed}床内 HGM}   & $\bigcirc$ \\
         \bottomrule
       \end{tabular}}\par
      \vspace{0.25em}
      {\tiny\raggedright $\times$ 移動床は失活追従不可、$\triangle$ 径方向流は機構複雑\par
       \vspace{0.18em}
       $\times$ 段間再加熱は床温を戻すのみで分単位の急速失活に間に合わない。切替再生で蓄えた熱を HGM が床内補償\par}
    \end{column}

    % ===== 右：反応器の特徴（上・小さめ）＋模式図（下）============
    \begin{column}{0.60\textwidth}
      % 特徴文：左のフロー図から離すため右へインデント、文字は小さめ
      \hspace*{9mm}\begin{minipage}{\dimexpr\linewidth-9mm\relax}
        \raggedright
        {\bfseries 反応器の特徴}\par
        \vspace{0.3em}
        {\scriptsize\raggedright\linespread{1.18}\selectfont
          ・\textbf{Catofin 型}の浅床軸流・多基切替式固定床を採用。\par
          ・吸熱反応と再生（コーク燃焼）を周期的に切替え、\textbf{計 14 基}で連続運転。\par
          ・床内の蓄熱体 \textbf{HGM} が反応熱を供給し床温を維持。\par
          ・大断面・低流速で接触時間を確保（$\tau\propto 1/u^{2}$、流速半減で 4 倍）し、短周期切替で失活に追従。\par
        }
      \end{minipage}\par

      \vspace{1.2em}   % 模式図を下げ、最下部をスライド下端に合わせる
      \centering
      {\bfseries\color{SpRed}浅床軸流・多基切替式固定床の模式図}\par
      \vspace{0.15em}
      \resizebox{0.95\linewidth}{!}{%
      \begin{tikzpicture}[
        >=Stealth,
        proc/.style={-{Stealth[length=2.4mm]}, line width=1.2pt},
        wall/.style={line width=1.3pt},
        lead/.style={line width=0.4pt, black!55},
        lbl/.style={font=\large, align=center},
      ]
        \foreach \dx/\toplab/\botlab/\stt/\bcol/\acol in {%
           0/{原料ガス}/{製品}/{反応中}/{ArrowBlue}/{ArrowBlue},
           5.4/{空気}/{燃焼ガス}/{再生中}/{orange!85!red}/{orange!85!red}%
        }{
          \begin{scope}[xshift=\dx cm]
            \def\R{1.72}\def\Hb{1.48}\def\ry{0.80}
            \draw[wall] (-\R,0)--(-\R,\Hb); \draw[wall] (\R,0)--(\R,\Hb);
            \draw[wall] (-\R,\Hb) arc[start angle=180, end angle=0, x radius=\R, y radius=\ry];
            \draw[wall] (\R,0) arc[start angle=0, end angle=-180, x radius=\R, y radius=\ry];
            \draw[line width=0.7pt] (-\R+0.07,\Hb-0.30)--(\R-0.07,\Hb-0.30);
            \fill[pattern=dots, pattern color=\bcol] (-\R+0.07,\Hb-0.80) rectangle (\R-0.07,\Hb-0.38);
            \draw[black!40] (-\R+0.07,\Hb-0.80) rectangle (\R-0.07,\Hb-0.38);
            \draw[line width=0.7pt] (-\R+0.07,\Hb-0.86)--(\R-0.07,\Hb-0.86);
            \draw[wall] (-0.17,\Hb+\ry-0.02)--(-0.17,\Hb+\ry+0.30);
            \draw[wall] (0.17,\Hb+\ry-0.02)--(0.17,\Hb+\ry+0.30);
            \draw[proc, \acol] (0,\Hb+\ry+0.92)--(0,\Hb+\ry+0.36);
            \node[lbl, above] at (0,\Hb+\ry+0.90) {\toplab};
            \draw[wall] (-0.17,-\ry+0.02)--(-0.17,-\ry-0.30);
            \draw[wall] (0.17,-\ry+0.02)--(0.17,-\ry-0.30);
            \draw[proc, \acol] (0,-\ry-0.36)--(0,-\ry-0.92);
            \node[lbl, below] at (0,-\ry-0.90) {\botlab};
            \node[font=\Large\bfseries, text=\acol] at (0,-\ry-1.74) {\stt};
          \end{scope}
        }
        \node[lbl, anchor=east, font=\scriptsize] (bl) at (-2.15,1.30) {触媒・HGM\\浅床};
        \draw[lead] (bl.east) -- (-1.60,1.00);
        \node[lbl, anchor=east, font=\scriptsize] (gl) at (-2.15,0.40) {支持グリッド};
        \draw[lead] (gl.east) -- (-1.60,0.62);
        \draw[{Stealth[length=2.6mm]}-{Stealth[length=2.6mm]}, line width=1.6pt, ArrowBlue]
              (1.95,0.95) -- (3.45,0.95);
        \node[lbl, ArrowBlue, font=\normalsize\bfseries, align=center] at (2.7,1.55) {周期\\切替};
        \node[font=\large, align=center, text=black] at (2.7,-3.45)
              {反応 $\rightleftarrows$ 再生 を周期交代　計 14 基};
      \end{tikzpicture}}
    \end{column}
  \end{columns}

\end{frame}
